% 原文 pp.91–98（印刷 pp.1209–1216）
\section{Lecture 7：二维玻色规范理论、Wilson 线算符与禁闭}

\begin{keytakeaways}
本讲把前两讲的二维规范理论从费米子情形转到玻色子情形：以 $U(1)$ 规范场耦合 $N$ 个复标量 $\phi_1,\dots,\phi_N$（线丛 $L$ 的截面）的拉格朗日量 (7.1) 为出发点，先按 $m^2$ 的两种极端讨论红外行为——$m^2\to+\infty$ 时四次项可忽略，模型回到第 5 讲所研究的有质量带电玻色子，它们只通过规范场弱耦合，玻色子间的一维 Coulomb 势导致禁闭，理论中所有态总电荷为零；$m^2\to-\infty$ 时 $|\phi|$ 被钉在 $a$ 上，低能有效理论是映入 $\mathbb{CP}^{N-1}$ 的 $\sigma$ 模型，而 $N=1$ 时 $\mathbb{CP}^{0}$ 只是一点，故理论有质量间隙与唯一真空。核心数学工具是：把 $\theta$ 依赖的 Fourier 系数实现为 $c_1\neq0$ 的线丛上的路径积分，用常微分方程边值问题论证瞬子的存在性、唯一性与确定尺度，从而得到与尺度无关的正作用量 $I$，再把它视为稀薄气体（dilute gas）求和，使配分函数化为对瞬子数 $m$ 与反瞬子数 $n$ 的求和（带权重 $W^{m+n}/(m!n!)$ 与 $\theta$ 相位）的指数化，得 $Z=e^{VP_0+2WP\cos\theta}$ 与 $E_{\mathrm{vac}}=-P_0-2Pe^{-I}\cos\theta$。由此得到本讲的主要结论：$\theta$ 依赖是 $e^{-I}\sim e^{-\mathrm{const}/e^2}$ 阶的非微扰效应，微扰论在任何有限阶都看不到 $\theta$；瞬子气体的稀薄与瞬子尺度的确定性使 $E_{\mathrm{vac}}(\theta)$ 光滑，故 $m^2\ll-e^2$ 时任何 $\theta$ 下都有质量间隙与唯一真空，而 $m^2\gg e^2$ 时 $\theta=\pm\pi$ 处导数不连续（有「割」），这个割必终止于某个 $m^2=m_c^2$，那里理论失去质量间隙。本讲最后引入 Wilson 线算符 $W_R(C)=\operatorname{Tr}|_R\operatorname{Hol}(A,C)$：它经典地规范不变，按 $A$ 的幂展开后逐项为多项式故可重整化，且只作乘性重整化，反正规化因子为 $e^{c'L(C)}$，在临界维数 4 中为 $e^{L(C)\Lambda f(e^2)+o(\Lambda)}$；在 Hamilton 图像下，位于 $t=0$ 切片内的圈只是「乘以一个函数」，而一般的圈须用带两个标记点 $P,Q$ 的 Hilbert 空间 $\mathcal{H}_{P,Q}$ 与 Hamilton 量 $H_{P,Q}$ 描述。
\end{keytakeaways}

% ===== 原文 p.91 =====
\subsection{二维中含玻色子的 $U(1)$ 规范理论的红外行为}
% 原文 §7.1：Infrared behavior of U(1) gauge theories with bosons in two dimensions

% 原文标题：Two-Dimensional Gauge Theory of bosons, the Wilson Line Operator, and Confinement
% （原书标题中 "bosons" 作小写，系原文排印如此，此处不作改动。）

上一讲我们考虑了带费米子的二维规范理论。今天我们将回到带玻色子的二维规范理论——关于后者，我们在第 5 讲已经有所了解。与以往一样，我们将在欧氏拉格朗日量（Euclidean Lagrangian）的框架下工作。

我们考虑二维中含玻色子的 $U(1)$ 规范理论。

我们的时空是一个黎曼面（Riemann surface）$\Sigma$。场为：$A$——某个线丛（line bundle）$L$ 上的一个 $U(1)$ 联络（connection）；以及 $\phi_1,\dots,\phi_N$——复标量场，它们是 $L$ 的截面（section）。

我们的拉格朗日量为
\begin{equation}
  \tag{7.1}
  \label{eq:7-1}
  \int d^2x\left(\frac{|{*}F|^2}{4e^2}+\frac12|d_A\phi|^2+\frac{\lambda}{4}\bigl(|\phi|^2-a^2\bigr)^2\right)-\frac{i\theta}{2\pi}\int F,
\end{equation}
其中 $|\phi|^2=\sum_i|\phi_i|^2$，$|d_A\phi|^2=\sum_i|d_A\phi_i|^2$。

我们要理解这个理论的红外行为。

经典地，把 $\phi_i$ 量子化所得的玻色子的「质量」是 $m^2=-\lambda a^2$。理论的参数是 $m,\theta,\lambda,e$（其中 $\lambda>0$）。下面我们研究系统在不同参数范围内的行为。

首先考虑 $m^2\to+\infty$ 而 $\lambda,e$ 固定的情形。此时拉格朗日量中的四次项变得无关紧要，模型本质上归结为第 5 讲研究过的模型：有质量的带电玻色子只通过规范场相互作用，而它们与规范场的耦合是弱的。理论有质量间隙（mass gap），玻色子之间的相互作用近似为一维 Coulomb 势（one-dimensional Coulomb potential），后者导致它们的禁闭（confinement）：理论中所有态的总电荷为零。此外，空间反转（space reversal）对称性在 $\theta=0$ 时保持，在 $\theta=\pi$ 时破缺，从而给出两个真空。作为 $m$ 与 $\theta$ 的函数，真空在 $\theta=\pi$ 处有一个「割」（cut）或不连续性。例如，真空能量密度（vacuum energy density）是 $\theta$ 的 $2\pi$ 周期函数，在 $-\pi\le\theta\le\pi$ 上由下面的公式 (7.2) 给出，而它在 $\theta=\pi$ 处的导数不连续。\footnote{本讲笔记由 Pavel Etingof 与 David Kazhdan 整理。}

% [Tutor Note: (7.1) 的 $\theta$ 项写作 $-\frac{i\theta}{2\pi}\int F$，而非四维常见的 $\frac{\theta}{8\pi^2}\int F\wedge F$。二维中 $F$ 是 2 形式，$\int_\Sigma F=2\pi c_1(L)$ 恰为 $2\pi$ 的整数倍，故该项在闭曲面上的贡献是 $-i\theta c_1$；路径积分权重 $e^{-S}$ 因而给出 $e^{i\theta c_1}$——这正是 p.1214 的 (7.4) 中每个瞬子带一个 $e^{i\theta}$ 的来源。若照搬四维归一化 $\frac{\theta}{8\pi^2}\int F\wedge F$，会得到 $e^{i\theta c_1/2}$ 之类的错误权重，并使 $\theta$ 的 $2\pi$ 周期性失效。另注意动能项写作 $\frac{|{*}F|^2}{4e^2}$：二维中 ${*}F$ 是 0 形式（标量），$|{*}F|^2$ 就是 $F$ 的模方，系数 $\frac{1}{4e^2}$ 只是约定（不少作者写 $\frac{1}{2e^2}$），但它与 (7.2) 的系数是绑定的，改写时不能只动一处。]

% [Tutor Note: 「$\theta=\pi$ 时空间反转破缺、产生两个真空」的理由值得点明：$E_{\mathrm{vac}}(\theta)$ 既是 $2\pi$ 周期的，又在 $\theta\to-\theta$ 下不变（空间反转把 $\theta$ 变号），故 $\theta=\pi$ 是这两个 $\mathbb{Z}_2$ 作用的公共不动点，左右两支在那里可以交叉而不必光滑相接；两支交叉即两个简并真空，空间反转在两者之间交换，于是自发破缺。这也说明作者为何强调 (7.2) 只是一个「$2\pi$ 周期函数在 $-\pi\le\theta\le\pi$ 上的取值」：抛物线 $\frac{e^2}{2}\bigl(\frac{\theta}{2\pi}\bigr)^2$ 本身并不周期，须按 $\theta\to\theta+2\pi$ 作周期延拓，延拓后在 $\theta=\pm\pi$ 处成为尖点，$\partial_\theta E_{\mathrm{vac}}$ 两侧分别为 $\pm\frac{e^2}{4\pi}$，导数跃变为 $\frac{e^2}{2\pi}$。]

% ===== 原文 p.92 =====
现在考虑 $m^2\to-\infty$ 的情形（同样固定 $\lambda,e$）。此时情况完全不同：由于四次项的存在，玻色子的经典势在 $|\phi|^2=a^2$ 处取极小，故极小值的空间是 $S^{2N-1}$。规范群 $U(1)$ 自由地作用在这个空间上，故经典基态的空间是 $\mathbb{CP}^{N-1}=S^{2N-1}/U(1)$。于是，低能有效理论是由 $\Sigma$ 到带圆度规（round metric）的 $\mathbb{CP}^{N-1}$ 的映射的理论（即 $\sigma$ 模型）。

\textbf{注（Remark）.} 在第 4 讲中我们看到，拉格朗日密度为 $\frac12(d\phi)^2+\lambda(\phi^2-a^2)^2$ 的理论在红外流向映入球面的 $\sigma$ 模型。我们曾指出，这意味着这个（超可重整的，superrenormalizable）理论是映入球面的 $\sigma$ 模型的一个良好截断。类似地，由拉格朗日量 (7.1) 定义的那个超可重整（因而严格存在）的理论，是映入 $\mathbb{CP}^{N-1}$ 的理论的一个良好的紫外截断。

我们已在第 5 讲研究了大 $N$ 的 $\mathbb{CP}^{N-1}$ 模型；我们发现，它在红外表现为 $N$ 个有质量玻色子的规范理论。今天我们将更详细地考察相反的极端，实际上就是 $N=1$ 的情形。此时 $\mathbb{CP}^{N-1}$ 是一个点，故由 (7.1) 定义的理论有质量间隙和唯一真空。问题在于理解该真空中物理的性质。尽管对我们来说，某些寻常问题在 $N=1$ 时是平凡的，我们仍会遇到全新的问题与方法。

\subsection{真空能量密度}
% 原文 §7.2：The vacuum energy density

回忆在第 5 讲中，对 $m^2\gg e^2$ 我们得到真空能量密度的如下取值：
\begin{equation}
  \tag{7.2}
  \label{eq:7-2}
  E_{\mathrm{vac}}(\theta)=\frac{e^2}{2}\left(\frac{\theta}{2\pi}\right)^2 .
\end{equation}

于是，能量密度对 $\theta$ 的导数正比于耦合 $e^2$ 的一次幂。

在 $m^2\ll-e^2$ 的情形，情况完全不同。即：现在路径积分中来自（相对于无穷远的）非平凡线丛的贡献在 $e^2$ 中是指数小的，因为这种丛的任何截面都必须在某处为零；而这个零点意味着，由于四次势的存在，作用量很大。因此，在微扰论的任何有限阶（按 $e^2$ 展开），路径积分、从而真空能量密度，都与 $\theta$ 无关。更确切地说，$\theta$ 依赖是指数小的，它来自对瞬子（instanton）的求和——非平凡第一 Chern 类（first Chern class）的最低作用量位形。

% [Tutor Note: 作者用「$m^2\gg e^2$」与「$m^2\ll -e^2$」划分参数区域，是把质量平方与耦合平方相比（二维中 $e$ 与 $a$ 的量纲均为质量，$m^2=-\lambda a^2$ 的量纲为质量平方），而不是与 $0$ 相比。两个极限的物理差别在于：前者 $m^2>0$ 使 $\phi$ 的径向模式很重、四次项可忽略，$\phi$ 只有微弱的长程规范相互作用；后者 $m^2<0$ 且 $|m^2|\gg e^2$ 使 $|\phi|$ 被牢牢钉在 $a$ 上，低能自由度只剩 $\mathbb{CP}^{N-1}$ 上的「Goldstone 型」模式（$N=1$ 时一个也没有），而径向激发有 $O(|m^2|)$ 的大质量。也正因如此，后文 p.1214 把同一条件简写作「$m^2\gg0$」与「$m^2\ll0$」在量纲上并不严格，应读回 $m^2\gg e^2$、$m^2\ll -e^2$。]

% [Tutor Note: 「$m^2\ll -e^2$ 时 $\theta$ 依赖只可能来自非平凡线丛」是纯拓扑的论证：$c_1(L)\neq0$ 的线丛上没有处处非零的光滑截面，故 $\phi$ 必在某点取零；而在 $|\phi|=0$ 处四次势取 $\frac{\lambda}{4}a^4$ 的大值，作用量因而很大。注意这里比较的是「截面的零点」与「势极小处 $|\phi|=a$」：$|m^2|$ 越大（势越深）论证越强。这正是「拓扑项不能被微扰地看到」的标准例子，可与四维 $\theta$ 真空中的瞬子按 $e^{-8\pi^2/g^2}$ 压低相对照（见 [I-Perturbative QFT] 相应章节）；区别在于二维的瞬子尺度是固定的，四维 Yang–Mills 瞬子则有尺度模。]

\subsection{瞬子}
% 原文 §7.3：Instantons

我们来计算带拉格朗日量 (7.1) 的路径积分，在它的 Fourier 展开中 $e^{i\theta}$ 项的系数。这等价于取 $c_1=1$ 的线丛上的联络的路径积分。

照例，对路径积分的主要贡献来自贴近「作用量极小位形」的场位形。这样的位形 $(\phi,A)$ 称为瞬子。

% ===== 原文 p.93 =====
设 $(\phi,A)$ 为一个瞬子。为使它的作用量有限，必须有
\begin{equation}
  \tag{7.3}
  \label{eq:7-3}
  |\phi(x)|\to a,\qquad F_A(x)\to 0\qquad\text{当 }x\to\infty,
\end{equation}
其中 $F_A$ 是 $A$ 的曲率。此外，由于 $\phi$ 经典地有质量，这一收敛是指数快的。因此，我们问题中的瞬子必须是一个高度局域化的场位形。

在下面关于瞬子的进一步讨论中，我们假定 $\Sigma=\mathbb{R}^2$，且 $c_1$ 是相对于无穷远处的一个平凡化（trivialization）来度量的。此时我们预期瞬子是旋转对称的（相对于某个旋转中心 $x_0$）。不失一般性可以设 $x_0=0$，于是瞬子具有形式 $\phi=ae^{i\alpha}f(r)$、$A=g(r)\,d\alpha$，其中 $f(r),g(r)$ 是半径 $r$ 的某些函数（$r,\alpha$ 是 $\mathbb{R}^2$ 上的极坐标）。容易看出，此时必须有 $f(0)=0$、$f(\infty)=1$，以及 $g(0)=0$、$g(\infty)=1$。于是我们得到一个关于常微分方程的边值问题。通过考察相应的 ODE 可以证明，该问题有唯一解。因此，瞬子在平移之下是唯一的。特别地，它有确定的尺度，其作用量是一个确定的正常数 $I$。

\textbf{注（Remark）.} 在这方面，我们的瞬子与 $\mathbb{CP}^{N-1}$ 模型的瞬子不同：后者可以被任何共形自同构（conformal automorphism）变换，因而不具有确定的尺度。

通过量纲分析（对空间、对 $\phi$、对 $I$ 作重标度）可以证明，$I$ 具有形式 $e^{-2}m^2h(\lambda/e^2)$。因此，若令 $e\to0$ 而 $m^2$ 与 $\lambda/e^2$ 固定，则作用量是 $O(1/e^2)$ 阶的。在这一极限下，瞬子贡献是指数小的。

% [Tutor Note: 拟设 $\phi=ae^{i\alpha}f(r)$、$A=g(r)d\alpha$ 与 $c_1=1$ 相容，可由 $g(\infty)=1$ 直接验算：绕无穷远的大圆 $\oint A=\int_0^{2\pi}g(\infty)\,d\alpha=2\pi$，故 $\int_{\mathbb{R}^2}F=2\pi$，即 $c_1=1$，与「$c_1$ 相对于无穷远的平凡化来度量」一致。若误把 $g(\infty)=1$ 写成 $g(\infty)=\frac12$（某些把 $A$ 当作「半联络」的约定），则 $\int F=\pi$，瞬子的 $\theta$ 权重会变成 $e^{i\theta/2}$，与 $\theta$ 的 $2\pi$ 周期性冲突。另注意 $\phi$ 前的整体因子 $a$：它保证 $|\phi|\to a$ 与 (7.3) 相符，也是后文 $\phi_*$ 归一化中处处出现的 $1/a$ 的来源。]

% [Tutor Note: 「由考察相应的 ODE 可证存在唯一解」是作者对跳步的典型省略。实际论证是边值问题的打靶法（shooting）：把 $f,g$ 的方程化为一阶系统，用 $r=0$ 处的正则性（$f\sim r$、$g\sim r^2$ 型）与 $r=\infty$ 处的指数衰减作两侧匹配，并用极值原理（maximum principle）证明单调性，从而匹配点唯一。这一步等价于说该瞬子没有除平移以外的集体坐标——与紧随其后的 Remark 恰成对照：$\mathbb{CP}^{N-1}$ 模型的瞬子有共形模，故其作用量与尺度无关（详见 §5 及 [II-CFT] 相关讨论）。]

\subsection{瞬子气体}
% 原文 §7.4：Instanton gas

现在我们要理解如何计算瞬子对路径积分的贡献，以及这种计算在何时给出好的近似。

首先考虑 $c_1>1$ 的线丛。可以证明，如果 $e^2/\lambda$ 充分大，这样的拓扑类中没有瞬子（但若 $e^2/\lambda$ 充分小，则有一个）。我们将假定 $c_1>1$ 时没有瞬子，但我们所说的可以推广到存在瞬子的情形（那只是需要把额外的对象纳入我们将要描述的稀薄瞬子气体（dilute instanton gas）之中）。在 $c_1>1$ 没有瞬子的情形，$c_1=2$ 时作用量的极小化问题没有全局极小（下确界取不到）。然而，我们可以考虑近似瞬子，其作用量任意接近下确界。更确切地说，若取
\[
A(x)=A_*(x-x_1)+A_*(x-x_2),\qquad \phi(x)=\phi_*(x-x_1)\phi_*(x-x_2)/a,
\]
其中 $(A_*,\phi_*)$ 是以 $0$ 为中心的瞬子，则 $(A,\phi)$ 的作用量等于 $2I$ 加上一个 $O\bigl(e^{-c|x_1-x_2|}\bigr)$ 的修正，其中 $c$ 是正常数。因此，若不存在 $c_1=2$ 的瞬子，则 $c_1=2$ 的作用量下确界为 $2I$。

类似地，对任何 $n\in\mathbb{Z}_+$，$c_1=n$ 的作用量下确界为 $nI$。这可由下述场位形的作用量推出：
\[
\left(\sum_{i=1}^{n}A_*(x-x_i),\; a\prod_{i=1}^{n}\frac{\phi_*(x-x_i)}{a}\right)
\]

% ===== 原文 p.94 =====
在 $|x_i-x_j|$ 很大时，该作用量指数地接近 $nI$。

现在考虑 $c_1<0$ 的情形。容易看出 $(-A_*,\bar\phi_*)$ 就是 $c_1=-1$ 的瞬子。它称为反瞬子（anti-instanton）。于是，$n<0$ 的情形与 $n\in\mathbb{Z}_+$ 的情形对称：场位形
\[
\left(-\sum_{i=1}^{|n|}A_*(x-y_i),\; a\prod_{i=1}^{|n|}\frac{\bar\phi_*(x-y_i)}{a}\right)
\]
在 $|y_i-y_j|$ 很大时作用量指数地接近 $|n|I$。

更一般地，可以考虑场位形
\[
\left(\sum_{i=1}^{n}A_*(x-x_i)-\sum_{j=1}^{m}A_*(x-y_j),\; a\prod_{i=1}^{n}\frac{\phi_*(x-x_i)}{a}\prod_{j=1}^{m}\frac{\bar\phi_*(x-y_j)}{a}\right),
\]
其 $c_1=n-m$，当 $x_i,y_j$ 彼此相距很远时作用量指数地接近 $(n+m)I$。这样的位形大致由位于点 $x_i$ 的 $n$ 个瞬子与位于点 $y_j$ 的 $m$ 个反瞬子组成。

\textbf{注（Remark）.} 这样的场位形构成所谓的「瞬子气体」（instanton gas）。下周，在另一类例子中（$2+1$ 维的 Polyakov 模型），我们将遇到粒子间有长程 Coulomb 力的「瞬子气体」，类似于等离子体中的情形。在目前的例子中，瞬子之间的相互作用在长程是指数小的；这更像普通原子构成的气体，例如大气中的空气。它是近乎理想的气体；理想气体是力严格为零的情形。任何真实气体（氢、氧）在密度足够小时都表现为近乎理想的气体，因为此时粒子通常相距很远，而那里相互作用很小。在本模型中瞬子正是如此，因为瞬子作用量很大（且瞬子有确定的尺度）。我们马上就会看到如何估计瞬子气体的密度；当 $e$ 很小（$m^2$、$\lambda/e^2$ 固定）时它是指数小的。

% [Tutor Note: 上页与这里的乘积拟设 $\phi=\prod_i\phi_*(x-x_i)/a$ 中，每多一个瞬子就多除一个 $a$，这是因为每个 $\phi_*$ 在无穷远处趋于 $a$：$n$ 个因子相乘在无穷远处趋于 $a^n$，除以 $a^{n-1}$ 才回到 $a$；作者把它写成 $a\prod_i(\phi_*/a)$，正是同一归一化。若漏掉这个 $1/a$，$|\phi|$ 在无穷远处会趋于 $a^{n}$，破坏 (7.3) 的边界条件，四次势将给出正比于 $(\Lambda_{\text{红外}})^2a^2$ 的假发散。]

\subsection{对瞬子求和}
% 原文 §7.5：Summing over instantons

如我们已指出的，带拉格朗日量 (7.1) 的路径积分按 $e^2$ 的幂所作的微扰级数不含瞬子的贡献，因为它们是指数小的。因此，微扰论中完全没有对 $\theta$ 的依赖。让我们引入一个精细化的微扰级数，它将把瞬子考虑在内。为此，我们将在体积为 $V$ 的二维时空上工作，并引入一个新的微扰参数 $W=Ve^{-I}$，它将刻画与 $\theta$ 有关的项。在这一精细化的微扰展开中，除了在平凡解（$A=0$、$\phi=$ 常数）附近构造微扰论之外，我们还在瞬子气体的位形附近展开。

在通常的微扰论中，人们构造的配分函数具有形式 $e^{VP_0}$，其中 $P_0$ 是耦合的函数，可在微扰论中算出。在像本讲这样有质量间隙的理论中，对体积的指数依赖有一个特别简单的直观解释。考虑在位置空间而非如通常那样在

% ===== 原文 p.95 =====
动量空间中计算 Feynman 图。在位置空间计算一张 Feynman 图，意味着对顶点在空间中的位置作积分，并在位置空间中求传播子。这正是 Feynman 图由路径积分导出时最直接地出现的表示（尽管实际计算 Feynman 图时人们通常用动量空间）。

给定质量间隙，所有传播子在长距离处都指数地趋于零。因此，在位置空间中计算一张连通 Feynman 图时，除指数小的效应外，所有顶点都落在一个有界的空间区域内（尺度为 $O(1/m')$，其中 $m'$ 是最小的质量）。连通 Feynman 图的「质心」可以位于空间中的任何地方，对它的位置积分便为每一张连通图给出一个因子 $V$（修正项出现在例如质心非常靠近空间边界时）。按通常的办法，对（可能不连通的）图求和给出连通贡献的指数，故配分函数的行为是 $e^{VP_0}$，其中 $P_0$ 是连通图的贡献。

现在考虑存在一个瞬子的情形。我们仍在位置空间中计算图。若一张 Feynman 图中顶点的质心不在距瞬子 $O(1/m')$ 的距离之内，则瞬子的存在并不改变该图的计算。若质心在 $O(1/m')$ 的距离之内，则微扰展开与没有瞬子时完全不同。于是，在存在瞬子的情形，我们得到的微扰级数是「真空」结果（不含瞬子）乘上一个因子，我们称之为 $P_+$，它依赖于瞬子附近微扰论被修改的方式。因此，含单个瞬子的配分函数的行为是
\[
Z_1=e^{VP_0}e^{-I}VP_+e^{i\theta},
\]
其中 $P_0$ 与 $P_+$ 要用微扰论计算，$I$ 是瞬子作用量，乘在 $P_+$ 上的因子 $V$ 来自对瞬子位置的积分，而我们已经把 $\theta$ 依赖包含在内。注意 $W=Ve^{-I}$ 的出现。\footnote{瞬子存在时对微扰级数的修正，对一张连通的 Feynman 图而言是相加的，而对未必连通的图则是相乘的。在 $Z_1$ 的公式中，$P_+$ 通过计算图中靠近瞬子的那些连通分支得到，而乘性因子 $e^{VP_0}$ 则来自把不靠近瞬子的那些分支的贡献取指数。} 若 $W\ll1$，单瞬子贡献与「真空」贡献相比小得可以忽略。若有人问：「$V$ 要多大，才有可观的概率看到一个瞬子？」答案的量级是 $V\sim e^{I}$，于是 $W\sim1$，单瞬子贡献与不含瞬子时算出的配分函数同量级。因此，单位体积内的瞬子密度是 $O(e^{-I})$ 阶的，并在 $e\to0$ 时指数小。

类似地，含单个反瞬子的路径积分的行为是 $e^{VP_0}e^{-I}VP_-e^{-i\theta}$，其中 $P_-$ 反映反瞬子附近微扰展开的修改。在目前这个理论中，由于只有 $\theta$ 项破坏定向反转（reversal of orientation）下的不变性，故有 $P_+=P_-$。（一般情形下只有 $P_-=\bar P_+$。）但我们将暂时把 $P_+$ 与 $P_-$ 当作彼此分离的常数。

现在考虑 $m$ 个瞬子位于点 $x_i$、$n$ 个反瞬子位于点 $y_j$ 的位形。我们假定 $x_i$ 与 $y_j$ 彼此之间的距离都远大于 $1/m'$。这是一个合理的假定，因为我们已经看到

% ===== 原文 p.96 =====
瞬子的密度是指数小的。于是，真空路径积分对每个瞬子乘一个因子 $P_+$、对每个反瞬子乘一个因子 $P_-$；而 Feynman 图在空间中的定域性、以及瞬子与反瞬子彼此远离，意味着它们的贡献可以独立地计算。对位置积分之后，每个瞬子或反瞬子给出一个因子 $W=Ve^{-I}$，同时瞬子给出 $e^{i\theta}$ 因子、反瞬子给出 $e^{-i\theta}$ 因子。

于是，这样得到的配分函数的近似具有形式
\begin{equation}
  \tag{7.4}
  \label{eq:7-4}
  Z(e^2,\theta,W)=e^{VP_0}\sum_{n,m=0}^{\infty}\frac{W^{m+n}P_+^mP_-^n e^{i\theta(n-m)}}{m!n!},
\end{equation}
（我们除以 $m!n!$，是因为 $m$ 个瞬子彼此全同、$n$ 个反瞬子也彼此全同，故在对它们的位置积分时，我们并不能为每一个瞬子或反瞬子都得到一个独立的因子 $V$。）

用通常的微积分对 (7.4) 求和，得到
\begin{equation}
  \tag{7.5}
  \label{eq:7-5}
  Z(e^2,\theta,W)=e^{VP_0+W(P_+e^{i\theta}+P_-e^{-i\theta})}.
\end{equation}

令 $P_+=P_-=P$，得到
\begin{equation}
  \tag{7.6}
  \label{eq:7-6}
  Z(e^2,\theta,W)=e^{VP_0+2WP\cos\theta}.
\end{equation}

现在我们可以在这个近似下计算真空的能量密度。这可以由等式 $Z=e^{-VE_{\mathrm{vac}}}$ 做到，其中 $E_{\mathrm{vac}}$ 是真空能量密度。于是，
\begin{equation}
  \tag{7.7}
  \label{eq:7-7}
  E_{\mathrm{vac}}=-P_0-2Pe^{-I}\cos\theta.
\end{equation}

如我们所预期的，真空能量的 $\theta$ 依赖被指数压低，但我们已能算出这一依赖的主要项。来自我们所忽略的效应的修正——例如瞬子之间的「相互作用」（比如它们的间距并不总是远大于 $1/m'$ 这一事实）——对 $E_{\mathrm{vac}}$ 的 $\theta$ 依赖给出比我们刚刚算出的更小的贡献。

在多大程度上可以改进对 $E_{\mathrm{vac}}$ 的 $\theta$ 依赖的这一估计？原则上，把 $P$ 按 $e^2$ 的微扰级数算到任意所需精度并无困难。然而，要计算真空能量中例如正比于 $\cos(2\theta)$ 的项，就必须对 $c_1=2$ 作精确计算。如果存在一个 $c_1=2$、作用量为 $I'$ 且 $I'<2I$ 的瞬子，这并无困难：此时 $E_{\mathrm{vac}}$ 含有一项 $\cos(2\theta)e^{-I'}P'$，可以像上面那样分析。然而，如果不存在这样的 $c_1=2$ 瞬子，则 $E_{\mathrm{vac}}$ 含有大致为 $\cos(2\theta)e^{-2I}$ 的项，它相当难以计算，来自两个瞬子之间的相互作用，特别是来自它们的间距未必远大于 $1/m'$ 这一事实。

我们回忆，对 $m^2\gg0$ 我们曾近似有 $E_{\mathrm{vac}}=e^2\theta^2/8\pi^2$（$|\theta|\le\pi$），特别地 $E_{\mathrm{vac}}$ 曾是 $\theta$ 的非光滑函数（在 $\theta=\pm\pi$ 处有割）。相反，在 $m^2\ll0$ 的情形，函数 $E_{\mathrm{vac}}(\theta)=-P_0-2Pe^{-I}\cos\theta$ 是光滑的。

% [Tutor Note: (7.4) 的指数 $e^{i\theta(n-m)}$ 与前后文不自洽，疑为排印笔误，应读作 $e^{i\theta(m-n)}$。检验：令 $m=1,n=0$（单个瞬子），(7.4) 给出 $W P_+e^{-i\theta}$，而 p.1213 明确写出 $Z_1=e^{VP_0}e^{-I}VP_+e^{+i\theta}$，且正文说「瞬子给 $e^{i\theta}$、反瞬子给 $e^{-i\theta}$」；再看 (7.5)，它由 (7.4) 求和而得，其指数是 $W(P_+e^{+i\theta}+P_-e^{-i\theta})$，只有当 (7.4) 的指数为 $i\theta(m-n)$ 时才与逐项求和相符。按印出的 $n-m$，(7.5) 将是 $e^{W(P_+e^{-i\theta}+P_-e^{i\theta})}$。]

% [Tutor Note: 由 (7.6) 到 (7.7) 用到了 $W=Ve^{-I}$，即 $VP_0+2WP\cos\theta=V\bigl(P_0+2Pe^{-I}\cos\theta\bigr)$，再与 $Z=e^{-VE_{\mathrm{vac}}}$ 比较。这里最易错的是量纲：$W$ 无量纲，而 $P_\pm$ 的量纲是能量密度（二维中为质量平方），故 $W P_\pm$ 的量纲为 $V\cdot(\text{能量密度})$，指数无量纲——若把 $W$ 也当成有量纲的量而漏掉 $V$，就会写出 $E_{\mathrm{vac}}\propto -2PW\cos\theta$ 这类量纲错误的式子。另注意 $\cos\theta$ 项对 $e^2$ 是本质奇异的（$e^{-I}\sim e^{-\mathrm{const}/e^2}$），故它不出现在任何有限阶微扰论中，这正是「微扰论看不到 $\theta$」的定量含义。]

% ===== 原文 p.97 =====
（诚然，我们只是在领头近似下算出了 $E_{\mathrm{vac}}(\theta)$，但相互作用的微弱、瞬子气体的稀薄以及瞬子尺度的确定性，都会使精确函数与我们算出的近似函数一样光滑。）因此，对 $m^2\ll0$，$\theta=\pm\pi$ 处没有割，并且对任何 $\theta$ 都有质量间隙和唯一真空。

于是，$\theta=\pm\pi$ 处的割必终止于某点 $m^2=m_c^2\in\mathbb{R}$。在这一点，理论很可能可以用一个具有四次临界点的有效势（effective potential）来描述，从而理论没有质量间隙。这类问题曾在第 6 讲末尾讨论过；很可能，一个具有第 6 讲末尾所得定性性质的有效势——对某个适当的有效低能自由度而言（该自由度不会与我们拉格朗日量中的任何基本场重合）——支配着我们今天所研究的系统在临界点附近的行为。

\subsection{Wilson 线算符}
% 原文 §7.6：The Wilson line operator

在本节及随后几节中，我们将定义 Wilson 线算符，并试图理解它的物理性质与形式性质。

假定我们有一个时空 $M$ 上规范群为 $G$ 的规范理论。设 $A$ 记相应的规范场，取值于 $G$ 的李代数 $\mathfrak{g}$。设 $C$ 是 $M$ 中的一条闭圈，并设 $\operatorname{Hol}(A,C)\in G$ 是联络 $A$ 沿 $C$ 的和乐（holonomy）（它只在共轭意义下有定义）。设 $R$ 是 $G$ 的一个有限维表示。定义经典的 Wilson 线（或 Wilson 圈，Wilson loop）泛函为
\begin{equation}
  \tag{7.8}
  \label{eq:7-8}
  W_R(C)(A)=\operatorname{Tr}|_R\bigl(\operatorname{Hol}(A,C)\bigr).
\end{equation}

显然这个泛函是规范不变的。

若 $C$ 是一些互不相连的圈 $C_i$ 的并，并标以表示 $R=(R_i)$，则按定义 $W_R(C)=\prod W_{R_i}(C_i)$。

若 $G$ 是交换群、且 $C$ 可缩，则 $W_R(C)=e^{i\int_D F}$，其中 $F$ 是 $A$ 的曲率，$D$ 是以 $\partial D=C$ 为边界的圆盘。

一个重要的推广如下。设 $\widetilde{G}$ 是 $G$ 的单连通覆盖（simply-connected cover），$R$ 是 $\widetilde{G}$ 的一个表示。此时容易看出，以同样方式定义的 $W_R(C)$ 在 $C$ 围成一个圆盘 $D$ 时仍然是良定义的。（在交换情形，这可由上面 $W$ 的积分表示推出。）

现在我们要在量子理论中定义 $W_R(C)$ 的类似物。为此需要重整化经典泛函 $W_R(C)$。这可以通过把 $W_R(C)$ 按 $A$ 的幂展开（如同把指数函数展开为 Taylor 级数）来做：
\begin{equation}
  \tag{7.9}
  \label{eq:7-9}
  W_R(C)=\dim(R)+\int dl\,\operatorname{Tr}|_R\bigl(A(l)\bigr)+\frac12\int dl\,dl'\,\operatorname{Tr}|_R\bigl(A(l)A(l')\bigr)+\cdots,
\end{equation}
其中 $l$ 是 $C$ 上的某个参数，$A(l)$ 是形式 $A$ 在点 $l$ 处、沿 $C$ 的切向量 $\frac{d}{dl}$（相对于这个参数化）上的取值。展开的每一项都是多项式，可以照常重整化。

事实上，在大多数情形下可以证明算符 $W_R(C)$ 只有乘性重整化（multiplicative renormalization）。粗略地说，这来自这样一个事实：经典地，$W_R(C)$ 是微分方程 $x'=Ax$ 沿该圈的单值性（monodromy）的迹。

% [Tutor Note: (7.9) 右端的 $\frac12\int dl\,dl'\,\operatorname{Tr}|_R(A(l)A(l'))$ 只在 $G$ 交换（或把 $A$ 当作 $c$ 数）时才是字面上的对称积分。非交换情形下 $W_R(C)=\operatorname{Tr}|_R P\exp\oint_C A$ 的 Taylor 展开带阶跃函数 $\theta(l-l')$，对称化后会多出 $\frac12\int dl\,dl'\,\theta(l-l')\operatorname{Tr}|_R\bigl([A(l),A(l')]\bigr)$ 型项。作者写成对称形式，是因为他紧接着只用「$W_R(C)$ 是 $A$ 的多项式泛函」这一点来论证乘性重整化；但若照 (7.9) 逐项作微扰重整化，必须补回路径排序，否则会漏掉 $\operatorname{Tr}(A^3)$ 型（三角型）图的贡献。]

% [Tutor Note: (7.8) 中作者特意加括号说明 $\operatorname{Hol}(A,C)$ 「只在共轭意义下有定义」：和乐依赖起点与平凡化的选取，换起点即作共轭变换 $h\to ghg^{-1}$，而 $\operatorname{Tr}|_R$ 在共轭下不变，故 (7.8) 才良定义。但也因此，在把 $W_R(C)$ 当作算符使用（例如取多个圈的乘积、或取期望值）时必须注意：$C$ 反向给出 $\operatorname{Hol}^{-1}$，即共轭表示 $\bar R$；有自交或相邻的圈需要 framing 才能定义。本讲后面用到的都是无自交的单圈，故这些微妙之处暂不出现。]

% ===== 原文 p.98 =====
这个方程被重整化为 $x'=(A+c)x$，其中 $c$ 是理论中的一个算符，在与 $A$ 相同的对称性下不变，处于规范群的伴随表示（adjoint representation）之中，其维数与 $A$ 相同或比 $A$ 更低。如果除常数之外没有这样的算符，那么这个方程在重整化下将变为 $x'=(A+c)x$，其中 $c$ 是一个标量算符。这表明 $W_R(C)$ 将具有乘性重整化；要重整化它，必须引入截断，并把 $W_R(C)$ 乘以一个函数 $e^{c'L(C)}$，其中 $L(C)$ 是 $C$ 的长度（或周长），$c'$ 是截断与耦合的函数。

更确切地说，可以证明在临界维数 4 中，发散的重整化因子具有形式 $e^{L(C)\Lambda f(e^2)+o(\Lambda)}$，其中 $\Lambda$ 是截断，$f(e^2)$ 可以用微扰论计算（但依赖于截断的精确选取）。在超可重整的情形（维数小于 4），发散因子按 $\Lambda$ 的较低次幂增长，或者根本不出现。

思考 Wilson 圈算符的一个可能更物理的方式，是把它等同于 Hilbert 空间上的一个算符。为此我们将使用 Hamilton 图像。于是我们取时空为 $M=M_s\times\mathbb{R}$，其中 $M_s$ 是空间流形，$\mathbb{R}$ 代表「时间」。理论的 Hilbert 空间 $\mathcal{H}$ 则是 $\mathcal{A}/\widetilde{G}$ 上的函数空间，其中 $\mathcal{A}$ 是 $M_s$ 中所有规范场的空间，$\widetilde{G}$ 是 $M_s$ 上规范变换的群，故 $\mathcal{A}/\widetilde{G}$ 是 $M_s$ 上联络的规范等价类的空间。空间 $\mathcal{H}$ 具有形式 $\mathcal{H}=(\mathcal{H}_0)^{\widetilde{G}}$，其中 $\mathcal{H}_0$ 是 $\mathcal{A}$ 上的函数空间。（在对 Hilbert 空间的这一描述中，我们实际上考虑的是不带物质场的「纯规范理论」；物质场可以以相当明显的方式纳入这些陈述。）

现在设 $C$ 是 $M$ 中的一条圈。首先设圈 $C$ 位于子流形 $t=0$ 之中。此时 $W_R(C)$ 是 $\mathcal{A}/\widetilde{G}$ 上的一个函数，作为 Hilbert 空间上的算符，它应被理解为乘以该函数。

现在考虑 $C$ 不在时空的水平截面之中、而是一条一般曲线的情形。此时 $W_R(C)$ 不再是乘以一个函数。让我们考虑相反的情形：$C$ 与切片 $t=0$ 横截相交。事实上，设 $C$ 在 $t=0$ 附近由 $P\times\mathbb{R}-Q\times\mathbb{R}$ 构成，其中 $P$ 与 $Q$ 是 $M_s$ 中的两个点，负号记录了 $C$ 与 $M_s$ 相交的定向。在 $t=0$ 附近我们仍然有一个与时间无关的系统，但它不再只由 $M_s$ 上的规范场（以及可能的其他场）组成；$M_s$ 上还有两个被标记的点，即 $P$ 与 $Q$，由于路径积分中的因子 $W_R(C)$，这两点处的物理是不同的。我们可以引入一个 Hilbert 空间 $\mathcal{H}_{P,Q}$ 和一个 Hamilton 量 $H_{P,Q}$ 来支配这一情形。但它们并不是通常的 Hilbert 空间与 Hamilton 量。

% [Tutor Note: 这一段在原书 p.1216 有一处明显的行文重复：先说方程重整化为 $x'=(A+c)x$（$c$ 为伴随表示中维数不高于 $A$ 的算符），紧接着又说「若除常数外没有这样的算符，则方程在重整化下变为 $x'=(A+c)x$，其中 $c$ 是标量算符」——两句给出的是同一个公式，属编辑疏漏。合理的读法是：前者为一般情形，后者为「当伴随表示中不存在维数 $\le\dim A$ 的非平凡算符时 $c$ 退化为常数」的特别情形。在二维中 $\dim A=1$、$F$ 的维数为 2，故维数 $\le1$ 的伴随算符只有 $A$ 自己（可由场重定义吸收），于是 $c$ 确实是常数，$W_R(C)$ 的反正规化只给出周长因子 $e^{c'L(C)}$——这就是为什么二维的 Wilson 圈只有周长发散而无尖点（cusp）反常维数。]

% [Tutor Note: 因子 $e^{L(C)\Lambda f(e^2)+o(\Lambda)}$ 的量纲值得核对：$L(C)$ 的量纲为长度（质量$^{-1}$），$\Lambda$ 为质量，故 $L(C)\Lambda$ 无量纲，$f(e^2)$ 必须是无量纲的耦合函数（四维中单圈值为 $f\propto e^2/8\pi^2$ 量级）。这也说明为何指数上是「周长」而不是「面积」：Wilson 圈的紫外发散来自沿圈传播的规范场自身，属局域于圈的量，与圈所围的面积无关；反之，若在指数上写面积，就会与「圈收缩到一点时 $W_R\to\dim(R)$」这一物理要求冲突。至于 $3+1$ 维中带尖点的圈，除周长因子外还有与尖点角度有关的额外发散，本书后面讨论 Wilson 圈期望值时才会用到。]
